\DTMsavedate{mydate}{2022-06-30}
\maketitle{Computer Architecture and Operating Systems}{Electronic Engineering}{\DTMusedate{mydate}}

% ----------------------------------------------------- %
% COURSE INFO
\subsection*{Course Information}
\vspace{0ex}%
\noindent\parbox{0.5\textwidth}{%
    \noindent\begin{tabular}{@{}ll}
                 \textsf{Course:}          & 	Computer Architecture and Operating Systems  \\
                 \textsf{Course Code:}         & 	TEE 5133   \\
                 \textsf{Credits:}    & 	10
    \end{tabular}}
\vspace{2ex}\hrule\vspace{2ex}

% ----------------------------------------------------- %
% INSTRUCTOR INFO
\subsection*{Lecturer Information}
\noindent\parbox{0.5\textwidth}{%
    \noindent\begin{tabular}{@{}ll}
                 \textsf{Name:}         &    Nomakhosi Ndiweni         \\
%\textsf{Phone:}        &    \phone          \\
\textsf{Email:}        &    nomakhosi.ndiweni@nust.ac.zw \\
\textsf{Qualifications:}        &  PhD in Electronic Engineering, MEng in Electronic Engineering, (Hon) BEng in \\
& Electronic Engineering
    \end{tabular}}
\vspace{2ex}\hrule\vspace{2ex}

% ----------------------------------------------------- %
% COURSE SYNOPSIS
\subsection*{Course Synopsis}
The module examines the evolution of computer hardware for Von Neumann machines.
It also covers operating systems for single tasking, process scheduling for concurrent operation, inter-process communication, deadlock avoidance and memory management.
The module also delves into virtual memory, architectures for parallel processing and computer networking.
\vspace{2ex} \hrule\vspace{2ex}

% ----------------------------------------------------- %
% LEARNING OBJECTIVES
% \subsection*{Learning Objectives}
% By the end of the course, candidates should be able to:
% \begin{outline}
%    \1      Understand the concept of states as applied to Control Systems.
% \end{outline}
% \vspace{2ex}\hrule\vspace{2ex}

% ----------------------------------------------------- %
% COURSE TOPICS
\subsection*{Topics}
\begin{outline}
    \1 The Evolution of Computer Systems
    \1 Hardware Architectures for Von Neumann Machines
    \1 Operating Systems for Single-Tasking Computers
    \1 Process Scheduling in Concurrent Operating Systems
    \1 Inter-Process Communication
    \1 Deadlock Avoidance in Multiprogrammed Computers
    \1 Memory Management in Multiprogrammed Machines
    \1 Hardware Architectures for Parallel Processing
    \1 Computer Networking
\end{outline}
\vspace{2ex} \hrule\vspace{2ex}

% ----------------------------------------------------- %
% Students Assenssment
\subsection*{Assessment of Students}
\begin{outline}
    \1 Continuous assessment will consist of research and presentations, assignments and  written in-class tests.
    \1 Final assessment will consist of an examination at the end of the semester.
\end{outline}
\vspace{2ex}\hrule\vspace{2ex}

% ----------------------------------------------------- %
% Grade Evaluation
\subsection*{Course Evaluation and Grading Policies}
Grades are based on:
Continuous assessment (\qty{25}{\percent}),
Final Exam (\qty{75}{\percent})
\vspace{2ex}\hrule\vspace{2ex}
% ----------------------------------------------------- %
% EQUIPMENT
% https://sites.udel.edu/computing-purchases/personal-specs/
% \subsection*{Essential Equipment}

% \vspace{2ex}\hrule\vspace{2ex}

% Policies and Other information
 \subsection*{Policies and Other Information}
% Key Text or some Books
 \subsection*{Key Text}

 \begin{itemize}
   \item Operating Systems: Design and Implementation, 3rd Edition by A. Tanebaum, A. Woodhull
 \end{itemize}
\vspace{2ex}\hrule
\vspace{2ex}\hrule\vspace{2ex}